\section{Normalization of the factor \texorpdfstring{$8\pi G$}{8piG}}
  \label{sec:normalization}

  \subsection{Two independent normalizations}
    \label{subsec:two-norms}

    The factor $8\pi G$ in equation~\eqref{eq:einstein} is entirely determined by
    two independent normalization conventions:

    \paragraph{Geometric normalization.}
      The companion paper~\cite{Gravity20} establishes that the infrared-dominant local
      part of the renormalized spectral entropy variation contains the Einstein tensor
      with coefficient $1/(16\pi G_N)$:
    \begin{equation}
      \frac{\delta\Spi^{\mathrm ren}}{\delta g^{\mu\nu}}\bigg|_{\mathrm IR,\,local}
      = \frac{1}{16\pi G_N}\,G_{\mu\nu}.
      \label{eq:geom-norm}
      \end{equation}
      This coefficient arises from the $a_2$ pole subtraction in the zeta-regularized
      determinant and the induced-gravity identification $G_N \sim 16\pi^2\ell_\chi^2$.

    \paragraph{Matter normalization.}
      The projected stress tensor is defined via the standard variational identity
    \begin{equation}
      T^{(\Pi)}_{\mu\nu} = \frac{2}{\sqrt{g}}
      \frac{\delta W_\Pi}{\delta g^{\mu\nu}},
      \label{eq:matter-norm}
      \end{equation}
      which carries the universal factor $2$ by convention.
      Equivalently, $\delta W_\Pi = \half\int\sqrt{g}\,T^{(\Pi)}_{\mu\nu}\delta g^{\mu\nu}$.

\subsection{The factor $8\pi G$ as a product of normalizations}
  \label{subsec:8piG}

  \begin{proposition}[Normalization identity]
    \label{prop:normalization}
    The equilibrium condition $\delta\mathcal{F} = 0$ with $\beta_*^{-1} = 1$ gives
    \begin{equation}
      \frac{1}{16\pi G_N}\,G_{\mu\nu}
      = \half\,T^{(\Pi)}_{\mu\nu},
    \end{equation}
    which is equivalent to $G_{\mu\nu} = 8\pi G_N\,T^{(\Pi)}_{\mu\nu}$.
    The numerical factor $8\pi G_N$ arises as the unique product of two independent
    normalizations:
    \begin{equation}
      8\pi G_N
      = \frac{1}{(16\pi G_N)^{-1} \times \tfrac{1}{2}}
      = \frac{1}{(16\pi G_N)^{-1}} \cdot 2
      \cdot 1,
      \label{eq:8piG-factorization}
    \end{equation}
    where the factor $(16\pi G_N)^{-1}$ comes from the geometric normalization
    \eqref{eq:geom-norm} and the factor $\frac{1}{2}$ from the stress-tensor
    definition \eqref{eq:matter-norm}.
    More transparently: from $\frac{1}{16\pi G_N}G_{\mu\nu} = \frac{1}{2}T_{\mu\nu}$
    one reads off directly $G_{\mu\nu} = 8\pi G_N T_{\mu\nu}$.
    No other factor enters.
  \end{proposition}

\subsection{Status of $\beta_*^{-1} = 1$}
  \label{subsec:beta-star-status}

  Setting $\beta_*^{-1} = 1$ in the constrained functional~\eqref{eq:F-functional}
  is not a dynamical assumption.
  It is a normalization choice equivalent to measuring energy in units where the
  spectral entropy and the projected matter functional are coupled with unit strength.
  Concretely, it amounts to absorbing $\beta_*^{-1}$ into the definition of
  $T^{(\Pi)}_{\mu\nu}$ as the effective stress tensor of the admissible projected
  sector, so that $W_\Pi$ is defined with the appropriate normalization.

  Alternatively, one may retain $\beta_*^{-1}$ as a free parameter and interpret
  $\beta_*^{-1} \neq 1$ as a departure from canonical spectral equilibrium.
  In that case,~\eqref{eq:EL-pointwise} generalizes to
  $G_{\mu\nu} = 8\pi G_N\,\beta_*^{-1}\,T^{(\Pi)}_{\mu\nu}$,
  which describes an off-equilibrium spectral geometry with an effective renormalization
  of the matter coupling.

\subsection{Relation between $\beta_*^{-1}$ and the local field $\beta(x)^{-1}$}
  \label{subsec:beta-relation}

  The local field $\beta(x)^{-1} = \beta_0^{-1} - \frac{1}{6}R(x)$
  and the global Lagrange multiplier $\beta_*^{-1}$ play distinct roles.
  The local field encodes spatially varying curvature corrections to the spectral
  multiplier; the global multiplier controls the overall strength of the
  entropy-matter coupling in the equilibrium functional.

  Their connection arises when the equilibrium condition is imposed locally rather
  than globally.
  Requiring $\delta\mathcal{F}$ to vanish not only for global metric variations but
  also for local ones supported near a point $x$ replaces $\beta_*^{-1}$ by
  $\beta(x)^{-1}$ in~\eqref{eq:EL-pointwise}, yielding
  \begin{equation}
    G_{\mu\nu}(x)
    = 8\pi G_N\,\beta(x)^{-1}\,T^{(\Pi)}_{\mu\nu}(x) + O(R\ell_\chi^2).
    \label{eq:local-einstein}
  \end{equation}
  To display the structure of the curvature correction transparently, we
  choose \emph{spectral units} in which $\beta_0^{-1} = 1$.

  \begin{remark}[Gauge choice $\beta_0^{-1} = 1$ and the spectral scale $t_*$]
    \label{rem:gauge}
    Setting $\beta_0^{-1} = 1$ corresponds, via $\beta_0^{-1} = 2/t_*$, to the
    choice $t_* = 2$ in spectral units.
    This is a \emph{gauge choice for the admissibility scale}: it fixes the unit
    of spectral resolution but does not affect any physical observable.
    In particular, the induced Newton coupling $G_N \sim 16\pi^2 \ell_\chi^2$
    established in~\cite{Gravity20} is independent of $t_*$, since it arises from
    the $a_2$ pole residue, which is a UV quantity fixed by the renormalization scheme.

    In these units, the local field simplifies to
    \begin{equation}
      \beta(x)^{-1} = 1 - \frac{1}{6}R(x) + O(\nabla^2 R),
      \label{eq:beta-spectral-units}
    \end{equation}
    and the locally imposed equilibrium condition~\eqref{eq:local-einstein} becomes
    \begin{equation}
      G_{\mu\nu}(x)
      = 8\pi G_N\!\left(1 - \tfrac{1}{6}R(x)\right) T^{(\Pi)}_{\mu\nu}(x)
      + O(R\ell_\chi^2).
      \label{eq:local-einstein-units}
    \end{equation}
    At zeroth order in $R\ell_\chi^2$, this reduces to the standard Einstein equation.
    The correction $-\frac{1}{6}R(x)$ is of order $R\ell_\chi^2$ and is absorbed into
    the $O(R\ell_\chi^2)$ remainder in the weak-curvature regime.
    Changing the gauge (i.e.\ changing $t_*$) rescales $\beta_0^{-1}$ and the
    correction simultaneously, leaving the physical content of~\eqref{eq:local-einstein}
    unchanged.
  \end{remark}
